\documentclass[11pt]{article}
\usepackage{graphicx}
\setlength{\parindent}{0pt}
\usepackage{hyperref}
\usepackage{enumitem}
\usepackage[utf8]{inputenc}
\usepackage[T1]{fontenc}
\usepackage[spanish]{babel}
\usepackage[left=1.06cm,top=1.7cm,right=1.06cm,bottom=1.7cm]{geometry}

\begin{document}

\begin{center}
\textbf{\Large Sergio Sebastián Pezo Jiménez}\\
\hrulefill
\end{center}

\begin{center}
Lima, Perú \textbullet\ +51 969 672 749 \textbullet\ sergiopezoj@gmail.com \textbullet\ linkedin.com/in/sergio-pezo \textbullet\ github.com/thsergitox
\end{center}

\vspace{4pt}

\begin{center}
\textbf{Perfil Profesional}
\end{center}

Estudiante de cuarto año de Ciencia de la Computación, con experiencia en proyectos de software e inteligencia artificial, me caracterizo por mi capacidad de adaptación y de agencia, sinceramente considero que cada problema tiene una solución, y mi objetivo es encontrarla. Actualmente me encuentro en búsqueda de una oportunidad que me permita aplicar mis habilidades, y a la vez, seguir aprendiendo y creciendo profesionalmente.

\vspace{10pt}

\begin{center}
\textbf{Experiencia Profesional}
\end{center}

\textbf{CoFoundy - Consultora de Software} \hfill Remoto

\textbf{Software Engineer} \hfill Mar 2025 – Dic 2025 

\begin{itemize}[noitemsep]
\item Lideré el diseño e implementación del backend principal de una plataforma para automatización de pedidos en restaurantes.
\item Migré un modelo manual basado en WhatsApp a un sistema digital escalable utilizado por decenas de clientes.
\item Reduje el tiempo de procesamiento de pedidos de aproximadamente 30 minutos a tiempo casi instantáneo.
\item Diseñé arquitectura de microservicios con NestJS, incluyendo lógica de negocio, autenticación y manejo de datos.
\item Implementé integración completa frontend-backend y participé en el desarrollo de interfaces.
\item Desplegué el sistema en producción con CI/CD (GitHub Actions + Railway).
\item Desarrollo del core del sistema con impacto directo en operaciones reales del negocio.
\end{itemize}

\vspace{10pt}


\begin{center}
\textbf{Educación}
\end{center}

\textbf{Universidad Nacional de Ingeniería} \hfill Lima, Perú

Ciencia de la Computación - 8.º ciclo, Décimo Superior \hfill Abril 2022 – Actualidad
\begin{itemize}[noitemsep]
\item Ingresé en 2021 en el examen de Ingreso Escolar Nacional (IEN), desde quinto año de secudaria a los 16 años, preparándome solamente con YouTube y libros. Lo que refleja la disciplina y esfuerzo que me caracteriza.
\item \textbf{Cursos relevantes:} Inteligencia Artificial; Programación Concurrente y Distribuida; Análisis y Diseño de Algoritmos; Estructura de Datos; Desarrollo de Software.

\end{itemize}

\textbf{Universidade Estadual de Campinas (UNICAMP) - Intercambio Académico} \hfill Sao Paulo, Brasil

Instituto de Computação \hfill Julio 2025 – Nov 2025
\begin{itemize}[noitemsep]
\item Durante este periodo de intercambio, interactué con profesores y estudiantes de Brasil de pregrado y posgrado, lo que me permitió ampliar mi perspectiva y conocimientos, así como trabajar con personas interdisciplinarias.
\item \textbf{Cursos Llevados:} Sistemas Distribuidos;Internet de las Cosas; Verificación, Validación y Testeo de Software; Bioinformática.
\end{itemize}

\vspace{10pt}


\begin{center}
\textbf{Proyectos}
\end{center}

\textbf{Sistema IoT con Azure Cloud Services} \hfill Oct 2025 – Nov 2025

Azure IoT Hub, Azure Functions, MQTT, WebSockets, Next.js \\

Proyecto realizado durante el intercambio académico en la UNICAMP para la materia de Internet de las Cosas, en conjunto con dos amigos de posgrado, donde yo me encargué en lo particular del software y la arquitectura. En cambio, ellos se encargaron de la parte física y electrónica.

\begin{itemize}[noitemsep]

\item Diseñé una arquitectura IoT escalable usando Azure IoT Hub como broker MQTT.
\item Implementé soporte para múltiples dispositivos y clientes mediante jerarquía de tópicos MQTT.
\item Desarrollé funciones serverless para procesamiento y almacenamiento de datos en tiempo real.
\item Implementé dashboard en Next.js con dominio propio y visualización en tiempo real.
\item Construí pipeline completo desde sensor físico hasta visualización web.
\end{itemize}

\textbf{Sistema de Log Distribuido con Algoritmos de Consenso} \hfill Oct 2025 – Nov 2025
Python, GCP, Lamport, Bully


Proyecto realizado durante el intercambio para la materia de Sistemas Distribuidos, donde se trabajó con Google Cloud Platform, logrando implementar y evaluar algoritmos de consenso como Lamport y Bully, con métricas de tiempo de respuesta y throughput.

\begin{itemize}[noitemsep]
\item Implementé sistema distribuido con relojes lógicos de Lamport para ordenamiento de eventos.
\item Desarrollé algoritmo Bully para elección de líder.
\item Desplegué sistema en múltiples máquinas virtuales y evalué comportamiento distribuido en diferentes zonas geográficas.
\end{itemize}

\textbf{TutorAI} \hfill May 2025 – Jul 2025

FastAPI, ElevenLabs

Realizado para la materia de Interacción Humano Computador en el séptimo ciclo, es un tutor para el aprendizaje de inglés usando IA, desarrollado con FastAPI y ElevenLabs, realizando pruebas de usabilidad con estudiantes universitarios que usaron y evaluaron la aplicación.

\begin{itemize}[noitemsep]
\item Desarrollé asistente de IA adaptativo para aprendizaje de inglés usado por estudiantes universitarios.
\item Implementé generación de contenido y feedback personalizado con APIs de voz.
\end{itemize}

\textbf{Sistema Multiagente con CRDTs} \hfill Sep 2024 – Nov 2024



\textbf{Backend Engineer} - Redis, LangGraph, Yjs
\begin{itemize}[noitemsep]
\item Diseñé arquitectura para edición colaborativa en tiempo real usando CRDTs.
\item Implementé APIs de agentes inteligentes.
\end{itemize}

\textbf{Plataforma de Formularios con IA} \hfill Feb 2024

\textbf{Backend Engineer}
\begin{itemize}[noitemsep]
\item Desarrollé sistema de generación dinámica de formularios con autenticación y roles.
\end{itemize}

\vspace{10pt}

\begin{center}
\textbf{Hackathons \& Premios}
\end{center}

\textbf{NASA Space Apps Challenge Campinas 2025} \hfill Semifinalista
\begin{itemize}[noitemsep]
\item Top 16 de más de 150 equipos.
\item Lideré equipo internacional (Brasil, Bolivia, Perú).
\item Implementé procesamiento concurrente de imágenes geoespaciales y visualización web en tiempo real.
\end{itemize}

\textbf{Hackathon Centro Internacional de la Papa} \hfill Finalista Mundial
\begin{itemize}[noitemsep]
\item Desarrollo de sistema de predicción agrícola usando ML.
\item Implementación completa del backend y pipeline de datos.
\end{itemize}

\textbf{Hackathon IMCA} \hfill 1.er lugar
\begin{itemize}[noitemsep]
\item Desarrollo de solución basada en redes neuronales para optimización de riego.
\end{itemize}

\vspace{10pt}

\begin{center}
\textbf{Liderazgo \& Actividades}
\end{center}

\textbf{ACECOM - Asociación Científica Especializada en Computación} \hfill Lima, Perú

\begin{itemize}[noitemsep]
\item Desarrollo de automatizaciones activamente utilizadas (cron jobs, integraciones WhatsApp/Discord).
\item Mentor de nuevos miembros en desarrollo de software.
\item Participación en eventos académicos explicando conceptos de redes y arquitectura.
\end{itemize}

\vspace{10pt}

\begin{center}
\textbf{Habilidades}
\end{center}

\textbf{Lenguajes:} TypeScript, Python, Java, Go. Realmente podría adaptarme a cualquier lenguaje de programación.

\textbf{Backend:} NestJS, FastAPI, Spring Boot, Express.

\textbf{Frontend:} React, Next.js, Remix, Astro.

\textbf{Cloud \& DevOps:} AWS, Azure, GCP, Docker, CI/CD.

\textbf{Bases de datos:} PostgreSQL, MongoDB, Redis. Realmente podría adaptarme a cualquier base de datos.

\textbf{IA:} TensorFlow, LangGraph, Claude Code, Codex.

\textbf{Idiomas:} Español (Nativo), Inglés (Avanzado), Portugués (Avanzado) \\

Mi sistema operativo favorito es Arch Linux con Hyprland, todos mis proyectos son desarrollados en este entorno.

\vspace{10pt}

\begin{center}
\textbf{Referencias}
\end{center}

Disponibles a solicitud.

\end{document}

